\chapter{Introduction}
% TODO: intro paragraph



\section{Motivation}
% - Find recent examples of things considered cults which have different impacts
%   in different countries.
% - Or find use of "cult" as a derogatory term.
% - Find use of "cult" as a positive/neutral term.
% - "Cult" is a highly ambiguous, contested term -- define ambiguity, and in
%   which way the concept is fuzzy. Consider the Hyperobject (Timothy Morton)
%   framing. Wittgenstein and family resemblance. MIVILUDES is itself aware of
%   this difficulty in its own framing.

- Seems to be an unstable word.(more)
- There was a moment of global interest - almost paranoia - in the 90s because of the Jonestown massacre, l'ordre du tempte solaire. 
``We know about `cults' largely by what the media tells us, and their views have been overwhelmingly negative\autocite{vandrielPrintMediaCoverage1988, pfeiferPsychologicalFramingCults1992}'' \autocite[pp.1]{dawsonCultsNewReligious2009}
- quick history of legal apparatus
- Trying to build a legal apparatus, the MIVILUDES ended up focusing on sectarian drifts: external negative consequences
- Signaling system, where things being signaled fall outside of New Religious Movements (as called by scholars), but can apply to other domain (list Miviludes reports categories)
- For example a leader Casasnovas has been signaled extensively and ended up being trialed for "wrongful practice of medecine". 
- In a nutshell, if we follow MIVILUDES's understanding of cults, we have now left a paradigm of cults as strictly spiritual groups, but entered 
- Cults seem to be live entities that adapt to their culture environment, akin to the way structuralists perceive myth:  local adaptation: Raël, the scientologists rely on sci-fi, ufologists, or l'ordre du temple solaire in early 90s build their eschatology\autocite[pp.180] {mayerOurTerrestrialJourney1999} on the culturally growing ecological anxiety. 
- Structuralist approach to cults
- The semantic space approach is structuralist par excellence

But what would it mean to look at cult strcuturally, with semantic space? 
- since the word is unstable in time and space, i propose to compare the structures of semantic spaces around the expression "cult" through 3 different epistemologies



- why are 
- what current shape can these cults take? 

\section{Research Goal}
% What questions shall be investigated and answered in this thesis? What is
% the scope of this work? -- comparing different epistemologies of "cult"
% (legal/MIVILUDES, scholarly/NRM studies, lay prototypes) through conceptual
% spaces / semantic embeddings.


But what would it mean to look at cult strcuturally, with semantic space? 
- since the word is unstable in time and space, i propose to compare the structures of semantic spaces around the expression "cult" through 3 different epistemologies:
- Scholarly work, to get a state of the art 
- Prototype theory approach to get current, live structures






Hypothesis the meaning of cult goes beyond NRM as forecasted by the reports to the MIVILUDES, spirituality not as a main driver. 


\section{Outline of Work}
% What is the structure of this thesis?
